\documentclass[11pt,a4paper]{article}

\usepackage[utf8]{inputenc}
\usepackage[T1]{fontenc}
\usepackage[english]{babel}
\usepackage{amsmath,amssymb}
\usepackage[top=1.2cm,bottom=1.5cm,left=2cm,right=2cm]{geometry}
\usepackage{graphicx}
\usepackage{booktabs}
\usepackage{hyperref}
\hypersetup{colorlinks=true,linkcolor=black,citecolor=blue,urlcolor=blue}

\title{Question D.4 --- Decision Protocol and Sensitivity to Censoring}
\author{}
\date{}

\begin{document}
\maketitle

\noindent Figures from the campaign of 2\,000 runs, reproduced by
\texttt{resultats\_R2/scripts/chiffres\_QD.py} and \texttt{derives\_QD.py}.

\section{Status of the Criterion, and Which Rule Governs}

The protocol asks for a decision criterion written down before any correlation is
inspected, and that claim cannot honestly be made here. The pre-registered document of this
work fixes the criteria of a \emph{benchmark}, the tests A1, A2, B1, B2, C1 and D1 to D3
qualifying candidate replacements for $\tau_{\mathrm{ARF}}$ against an analytic reference.
It contains neither the indicator evaluated here, nor its comparators, nor the decision
domain $\Delta e \ge 0.10$, nor the threshold on $\tau_b$, nor the rule on interval bounds,
and the columns \texttt{dans\_domaine\_decision}, \texttt{seuil\_detectabilite} and
\texttt{discernable} come from the same analysis pass as the coefficients they qualify.
Anteriority is therefore not claimed. The criterion is fixed here, applied uniformly, and
no cell is reclassified afterwards.

Two rules coexist in this work: a \emph{threshold rule}, $|\tau_b| > 0.13295637$, and an
\emph{interval rule}, the bootstrap interval excludes zero. They disagree on three cells
out of $108$: at $\Delta e = 0.287$, $\tau_{\mathrm{swap}}(75\%)$ has $\tau_b = +0.1363$
with interval $[-0.0068 ; +0.2704]$, discernible under the threshold rule only, while
$S_{\max}(H)$ has $+0.1308$ with $[+0.0027 ; +0.2519]$, discernible under the interval rule
only; at $\Delta e = 0.391$, $\tau_{\mathrm{swap}}(75\%)$ has $+0.1240$ with
$[+0.0010 ; +0.2505]$, again interval only.

\textbf{The interval rule governs.} A fixed threshold on $|\tau_b|$ ignores the precision
of the estimate, treating a coefficient of $0.13$ with a wide interval like one measured
tightly. It is also a function of $n$ alone, so it says nothing about the tie structure of
the pair tested, and Q~D.2 established that this structure differs systematically between
the $\tau_{\mathrm{swap}}(q)$ comparators and the never-censored $A(H)$ and $S_{\max}(H)$.
The pre-registered benchmark already uses the interval rule for test B1, so adopting it
keeps one rule across the project. Under it, $A(H)$ is discernible at $0$ of $18$
amplitudes, unchanged since both rules agree on every $A(H)$ cell, and $S_{\max}(H)$ at
$\mathbf{15}$ of $18$; the \texttt{discernable} column implements the threshold rule and
would give $14$, differing on $\Delta e = 0.287$ alone. The count reported throughout Q~D
is $15$, and none of the three disputed cells is quoted individually as a result.

\paragraph{What ranking censored runs last rests on.} The horizon $H = 2\,000$ is identical
for every run, so a duration not reached genuinely exceeds every observed one and censored
runs may take the last rank, tied among themselves. The argument depends entirely on that:
under a variable horizon a censored value would only be known to exceed \emph{its own}
horizon, a run censored at $H = 500$ and an observed run of $1\,200$ steps could not be
ordered at all, and the correct tool would be a censoring-aware estimator such as Harrell's
$C$ index. The convention used here rests on a property of this design and does not
transfer.

\section{Censoring by Amplitude, and Whether It Drives Anything}

Aggregate rates conceal what matters, so both the grid and the decision domain are
reported.

\begin{table}[htbp]
  \centering
  \begin{tabular}{lccc}
    \toprule
    Quantity & overall & worst on grid (at $\Delta e$) & worst in domain (at $\Delta e$) \\
    \midrule
    $\tau_{\mathrm{ARF}}$, $A(H)$, $S_{\max}(H)$ & $0.0000$ & $0.0000$ & $0.0000$ \\
    $\tau_{\mathrm{swap}}(25\%)$ & $0.0005$ & $0.0100$ ($0.141$) & $0.0100$ ($0.141$) \\
    $\tau_{\mathrm{swap}}(50\%)$ & $0.0110$ & $0.1200$ ($0.085$) & $0.0600$ ($0.141$) \\
    $\tau_{\mathrm{swap}}(75\%)$ & $0.0910$ & $0.7800$ ($0.085$) & $0.5400$ ($0.141$) \\
    $\tau_{\mathrm{det}}$, $\lambda = 8$  & $0.0470$ & $0.9000$ ($0.028$) & $0.0000$ \\
    $\tau_{\mathrm{det}}$, $\lambda = 25$ & $0.6110$ & $1.0000$ & $1.0000$ \\
    $\tau_{\mathrm{det}}$, $\lambda = 50$ & $0.9995$ & $1.0000$ & $1.0000$ \\
    \bottomrule
  \end{tabular}
  \caption{The two right-hand columns differ because the two amplitudes below
  $\Delta e = 0.10$ carry most of the censoring and lie outside the domain on which the
  criterion is declared. Baseline over $3\,000$ pre-drift steps.}
  \label{tab:censure}
\end{table}

The benchmark criterion of this work requires censoring of at most $5\%$ over the declared
domain. Inside the decision domain $\tau_{\mathrm{swap}}(75\%)$ reaches $54\%$ at
$\Delta e = 0.141$, exceeding it by a factor of $10.8$; its worst grid value, $78\%$ at
$\Delta e = 0.085$, is quoted as a grid measurement since that amplitude is outside the
domain. Its coefficients appear in Q~D.3 because suppressing them would hide the comparator
that most stresses the method, but with the failure stated: at $\Delta e = 0.141$ its
$\tau_b$ rests on a minority of observed runs and is not quoted as evidence. For the same
reason no median detection delay is reported at $\lambda = 25$ or $\lambda = 50$, nor at
$\lambda = 8$ for $\Delta e = 0.028$ where censoring reaches $90\%$: a median over
surviving runs excludes the slowest by construction.

\paragraph{Does censoring drive the results? No.} Every coefficient was recomputed with
censored runs imputed to the last rank and again on complete cases only. For the three
never-censored comparators the two agree to $0.000000$, which they must, and which controls
the implementation. Elsewhere the largest disagreement in the whole table is $0.040259$ on
$\tau_{\mathrm{swap}}(75\%)$, then $0.030591$ on $\tau_{\mathrm{swap}}(50\%)$ and
$0.000228$ on $\tau_{\mathrm{swap}}(25\%)$: an order of magnitude below the coefficients
carrying the conclusions of Q~D.3 and well inside their intervals. The imputation
convention could be replaced by complete-case analysis without changing a verdict. One
correction to an earlier internal note: ``agree to three decimal places'' is too strong, as
agreement is to two decimal places for $\tau_{\mathrm{swap}}(50\%)$ and
$\tau_{\mathrm{swap}}(75\%)$.

\section{The Blind Spot, Quantified}

At $\lambda = 8$, $1\,906$ runs out of $2\,000$ alarm ($4.70\%$ silent); at $\lambda = 25$,
$778$ ($61.10\%$ silent); at $\lambda = 50$, one single run ($99.95\%$ silent). That is why
the campaign fixes no $\lambda$: with $99.95\%$ censoring no correlation involving
$\tau_{\mathrm{det}}$ is computable at that threshold, and a threshold frozen in the
simulation loop, as an earlier prototype did, would have produced a $\tau_{\mathrm{det}}$
column empty in all but one row.

Read as the race of question~A, on uncensored runs,
$\tau_{\mathrm{det}} < \tau_{\mathrm{ARF}}$ in $95.44\%$ of runs at $\lambda = 8$, $6.56\%$
at $\lambda = 25$, $0\%$ at $\lambda = 50$, so the transition from ``never late'' to
``never fires'' happens within a narrow band of the tuning parameter. The rate is not
stationary in amplitude, though: at $\lambda = 8$ it is $0\%$ at $\Delta e = 0.028$, $62\%$
at $\Delta e = 0.085$, and $88\%$ to $99\%$ above. The blind spot depends on the pair
$(\lambda, \Delta e)$, and lowering $\lambda$ does not remove it at the bottom of the grid.

\paragraph{Not addressed.} No result of Q~D holds as a function of $M$, the campaign
existing at $M = 10$ only, and none involves $\tau_{\mathrm{det}}$ at $\lambda = 25$ or
$\lambda = 50$ as a correlated variable, for the censoring reason above. Harrell's $C$
index was not used: the common-horizon argument makes rank imputation licit here and the
complete-case comparison shows the choice is immaterial, though under a variable horizon
that reasoning would fail and the $C$ index would be required. The censoring rates of
Table~\ref{tab:censure} depend on the baseline window, whose two variants are reported in
Q~D.1 and shift them by under one point.

\end{document}
